\documentclass[../../../include/open-logic-section]{subfiles}

\begin{document}
\olfileid{sth}{ordinals}{zfm}	
\olsection{$\ZFminus$：一个里程碑}

如何为替换公理提供证成（如果可以证成的话），这个问题并不简单。因此，我们将把这个问题留到第 \olref[replacement][]{chap} 章讨论。然而，随着替换公理的加入，我们达到了又一个重要的里程碑。我们现在已经具备了理论 $\ZFminus$ 所需的全部公理。具体来说：

\begin{defn}
理论 $\ZFminus$ 包含以下公理：外延公理、并集公理、对集公理、幂集公理、无穷公理，以及分离公理模式和替换公理模式的所有实例。换言之，$\ZFminus$ 就是在 $\Zminus$ 的基础上添加替换公理得到的。
\end{defn}

\noindent
$\ZFminus$ 表示 \emph{Zermelo（策梅洛）--Fraenkel（弗兰克尔）} 集合论（\emph{减去}我们稍后会讲到的东西）。Fraenkel 获此殊荣，是因为替换公理的表述归功于他 \citeyear{Fraenkel1922}，不过第一个精确的表述是由 \citet{Skolem1922} 给出的。

\end{document}